\section{Benchmark Configuration}

Benchmark inputs are stored as JSON files in:

\begin{lstlisting}
input_json_files/
\end{lstlisting}

The main Module A paper-like configuration is:

\begin{lstlisting}
input_json_files/benchmark_tao2022_mof5_co2.json
\end{lstlisting}

\subsection{Configuration Sections}

\begin{itemize}
    \item \texttt{material}: framework name, CIF source, charge scheme, and source dataset.
    \item \texttt{adsorbates}: adsorbate components and optional mixture definition.
    \item \texttt{conditions}: temperature and pressure grid.
    \item \texttt{benchmark}: task, metrics, paper metadata, and reference-data configuration.
    \item \texttt{simulation}: LAMMPS/GCMC settings.
    \item \texttt{evaluation}: absolute/excess handling and gas-density model.
    \item \texttt{output}: output directory, run ID, overwrite behavior, and saved artifacts.
\end{itemize}

\subsection{Evaluation Configuration}

The current evaluation block is:

\begin{lstlisting}
"evaluation": {
  "simulation_basis": "absolute",
  "report_excess": true,
  "pore_volume_cm3_g": "auto",
  "pore_volume_source": null,
  "pore_volume_method": null,
  "gas_density_backend": "HEOS",
  "reference_matching": "by_basis"
}
\end{lstlisting}

With \texttt{pore\_volume\_cm3\_g} set to \texttt{auto}, the run planner reads the accessible pore volume from CRAFTED for the resolved material.
For IRMOF-1, this currently resolves to:

\begin{equation}
    V_{\mathrm{pore}} = 0.845545~\mathrm{cm^3~g^{-1}}.
\end{equation}

\subsection{Output Directories}

Runs are written below:

\begin{lstlisting}
outputs/module_A_co2_isotherm/
\end{lstlisting}

To avoid overwriting older results, production runs should use either \texttt{--new-run} or a unique \texttt{--run-id}.
